% Section 6 (Conclusion) of "Conglomerate Mergers under Multiplicative Growth".
% 15 Nov 2026 draft from the front-matter doc (C1-C4 as ruled). Simon writes the final prose over this.
\section{Conclusion}
\label{sec:Conclusion}

This article has developed a foundation for conglomerate mergers grounded in multiplicative dynamics and time-average growth. Building on \textcite{peters2022ergodicity}, who show that pooling resources is positive-sum under multiplicative dynamics, we have embedded their result in an agent-based model of industry dynamics in which firms must agree to pool on a finite record, can be selected out at a floor, and pay a cost for coordination. In that setting growth-rate optimization through variance reduction explains conglomerate formation without economies of scale, market power, tax motives or managerial incentives, and acting on the incentive compresses the size distribution within markets: the tail exponent rises with the sharing rate and the within-market Herfindahl index falls in every shock family. Unlike portfolio diversification, which reduces a shareholder's exposure while leaving each firm's growth path unchanged, cooperative diversification through pooling raises the growth rate of the constituent firms themselves. This extends the risk-reduction arguments of \textcite{amihud1981risk}, \textcite{lewellen1971pure} and \textcite{mueller1969theory} to the framework appropriate for multiplicative processes, and complements the internal capital markets literature \parencite{stein1997internal, gertner1994internal} by identifying a growth-rate channel that operates without information asymmetries or financing frictions.

The model's results are about who pools with whom and what it does for them. Pooling moves arithmetic growth from volatile members to calm ones, so it is the volatile side that refuses: mergers between existing units fail on that veto, and units grow by recruiting small firms. For the firms that join, membership raises time-average growth and cuts the hazard of being recapitalized at the floor by an order of magnitude; half to four fifths of the gain is ruin that did not happen, which no comparison of typical periods can see --- a measured reason why cross-sectional evaluations find a conglomerate discount even where conglomeration pays, and why the discount reverses in a crisis. Units that set their own sharing rate choose about 0.9 and never one, because members of unequal volatility must agree; the rate is high because in the model, as in no real firm, sharing carries no incentive cost. The distributional footprint is a tail exponent that rises in proportion to the sharing rate, within-market concentration that falls, ranks that move less, and a top decile of owners that is somewhat larger and markedly more persistent --- persistence that is lower ruin hazard, not market power, since no unit holds two firms in one market. On scope, the verdict is this: conglomerates form at once and stay small, not because the environment stops them --- with consent switched off the same floor permits units larger than the equal-members reference under light-tailed shocks --- but because incumbents refuse the floor-sized recruits whose admission would pay the recruit and not them; attrition takes over as the binding constraint only under the heaviest tails.

Several extensions follow from what the model leaves out. A pricing mechanism would allow statements about consumer welfare, which a framework confined to market structure cannot make. Correlation that rises in a crisis, rather than the constant correlation tested here, is the case in which the insurance is most valuable and thinnest at once. An incentive cost of the sharing rate itself --- the socialism of internal capital markets --- would bring the chosen rate down from the insurance-only optimum toward what firms do, and would let the model speak to how deep reallocation goes, not only to whether it pays. The exclusion the results identify is a contracting problem, and an admission price or a member-specific sharing rate is the natural next contract to put in front of the same firms. On the empirical side, the prediction that does not depend on a stationary environment is that scope falls with environment volatility through attrition while the willingness to pool rises with it; both halves are observable in merger and exit data by sector. Finally, whether firms that judge mergers on records of different length can coexist, and which length survives, is a question the model is built to ask and this paper has not asked.
